% -----------------------------------------------------------------------------------
\section{Robustness of the within-problem conflict}\label{sec:robust}

The central conflict between \cri{1} and \cri{2} was obtained with deterministic policies, a single
always-reached choice, and a predictor targeting that choice. This section asks which of these
features it needs.

\subsection{Imperfect prediction}

By Proposition~\ref{prop:newcomb}, the conflict holds for every $0\le\varepsilon<\tfrac14$ at the
scaled payoffs, and for every $\varepsilon<\tfrac12$ at suitable stakes, since $m(1-2\varepsilon)>c$ can
be arranged whenever $\varepsilon<\tfrac12$. At $\varepsilon=\tfrac12$ the symmetric predictor carries no
policy dependence and this mechanism disappears. Theorem~\ref{thm:characterisation} gives the
general statement: what is needed is not accuracy but any difference between the state laws that
two prescriptions induce.

\subsection{Randomisation does not reconcile ex-ante and causal evaluation}

Let the agent one-box with probability $p$. Suppose first that the predictor observes the
distribution, filling the box with probability $\varepsilon+(1-2\varepsilon)p$, and that the action is
drawn independently given $p$. Expanding,
\[
  V(p)=p\cdot2\bigl(\varepsilon+(1-2\varepsilon)p\bigr)+(1-p)\bigl(2\varepsilon+2(1-2\varepsilon)p+1\bigr)
      =1+2\varepsilon+(1-4\varepsilon)p .
\]
The same formula holds if the predictor instead observes the sampled action and predicts it with
symmetric error $\varepsilon$. For $\varepsilon<\tfrac14$, \cri{1} requires $p=1$ and \cri{2} requires $p=0$.
The conflict also admits an exact quantitative form. Put $g=1-4\varepsilon>0$, let
$R_1(p)=\max_qV(q)-V(p)$ be ex-ante regret, and let $R_{\mathrm{causal}}(p)$ be the expected regret
relative to the best action in the realised state.

\begin{proposition}[Exact trade-off]\label{prop:tradeoff}
In randomised Newcomb, $R_1(p)=g(1-p)$ and $R_{\mathrm{causal}}(p)=p$, so that
\[
  \frac{R_1(p)}{g}+R_{\mathrm{causal}}(p)=1 .
\]
Simultaneous tolerances $R_1\le\delta_1$ and $R_{\mathrm{causal}}\le\delta_2$ are achievable if and only
if $\delta_1/g+\delta_2\ge1$. Under infallible prediction, every policy has regret at least $\tfrac12$
by one of the two measures.
\end{proposition}

\begin{proof}
$V$ is affine with slope $g$ and maximal at $p=1$, which gives $R_1$. In every realised state
two-boxing is better by exactly one unit, so the causal regret is $1$ when the agent one-boxes and
$0$ otherwise. The tolerances require $1-\delta_1/g\le p\le\delta_2$, a nonempty interval exactly
under the stated inequality. At $\varepsilon=0$, $g=1$ and $\max(R_1,R_{\mathrm{causal}})\ge\tfrac12$.
\end{proof}

\subsection{Existence of optimal and ratifiable mixed policies}\label{sec:mixedexist}

Admitting mixed policies does not by itself make either \cri{1} or \cri{3} satisfiable; regularity
of the predictor's dependence on the distribution is what matters.

\begin{proposition}\label{prop:mixedfail}
\begin{enumerate}[label=(\roman*)]
  \item Let a binary reward, independent of the realised action, succeed with probability $q(p)=p$ for
        $p<1$ and $q(1)=0$. Then $V(p)=q(p)$ has supremum $1$ and no maximiser.
  \item Let Death choose city $1$ if and only if $p\ge\tfrac12$, where $p$ is the probability of city $1$,
        with payoff $1$ for avoiding Death. No mixture is ratifiable, although Death observes only the
        distribution.
\end{enumerate}
\end{proposition}

\begin{proof}
(i) is immediate. (ii) If $p<\tfrac12$, Death is in city $0$, and city $1$ is the unique causal best
response, so a ratifiable $p$ would put all its mass on city $1$; if $p\ge\tfrac12$, city $0$ is the
unique best response. Either way $p$ is not supported on best responses to itself.
\end{proof}

\begin{theorem}[Existence under continuity]\label{thm:mixedexist}
Let actions and states be finite, let the action be sampled independently given the mixed policy
$p\in\Delta(A)$, and let $p\mapsto\mu_p\in\Delta(S)$ be continuous. Then a mixed policy satisfying
\cri{1} exists, and so does a mixed policy all of whose supported actions are causal best responses
to $\mu_p$.
\end{theorem}

\begin{proof}
$V(p)=\sum_ap_a\sum_s\mu_p(s)u(s,a)$ is continuous on the compact simplex, so it attains its maximum.
For ratifiability, let $Q(a;p)=\sum_s\mu_p(s)u(s,a)$ and $B(p)=\Delta\bigl(\argmax_aQ(a;p)\bigr)$. The
correspondence $B$ has nonempty, compact, convex values, and its graph is closed because the
finitely many functions $Q(a;\cdot)$ are continuous. By Kakutani's theorem
\autocite{kakutani1941} there is $p\in B(p)$.
\end{proof}

Theorem~\ref{thm:mixedexist} is the one-shot core of two published existence results with
stronger structure: Halpern's theorem on ratifiable acts, which assumes costless statewise mixing
and conditional state laws that are affine in the mixture \autocite{halpern2020}, and the theorem of
\textcite{belletal2021} that every continuous Newcomblike decision process has a strongly ratifiable
policy. Neither identifies ratifiable policies with ex-ante maximisers, and
Proposition~\ref{prop:tradeoff} shows that they cannot be identified.

\begin{remark}[Predicting a realised sample]\label{rem:samples}
If Death independently samples the agent's distribution, the usual symmetric model admits
$p=\tfrac12$. If Death instead predicts the agent's realised private sample, the verdict depends on
further semantics: whether the agent knows its sample, and whether a deviation preserves or
replaces the randomness the predictor tracks. Conditioning on a revealed, perfectly predicted
sample restores the deterministic instability. A marginal test of support-optimality can still
accept $p=\tfrac12$ while ignoring the correlation between the prediction and the sample, and should
not be confused with a criterion that preserves that correlation when the unchanged policy is
evaluated. Skyrms already cautioned against an unconditional randomisation rescue
\autocite{skyrms1990ratif}.
\end{remark}

\subsection{Restrictions on what predictors may target}\label{sec:targets}

Some responses to Newcomb-type problems restrict prediction to information sets that are in some
sense actual. Newcomb's problem has one information set, reached with probability one under every
policy, so its conflict survives every such restriction except the last:

\begin{center}
\small
\begin{tabular}{@{}>{\raggedright\arraybackslash}p{19em}c@{}}
  \toprule
  admissible prediction targets & Newcomb conflict survives \\
  \midrule
  information sets positively reached under the evaluated policy & yes \\
  information sets reached under every admissible policy & yes \\
  information sets on the realised path & yes \\
  information sets already reached before the prediction occurs & no \\
  \bottomrule
\end{tabular}
\end{center}

The last restriction removes anticipatory dependence altogether: a \enquote{prediction} of an action
already taken is an ordinary causal response, which belongs in the structural plan, so $\mu_\pi$ no
longer depends on $\pi$. In that case ex-ante and causal evaluation agree.

\begin{proposition}\label{prop:causaltree}
If $\mu_\pi=\mu$ does not depend on $\pi$, then every \cri{1}-optimal policy satisfies \cri{3},
and therefore \cri{2}.
\end{proposition}

\begin{proof}
By Lemma~\ref{lem:local},
\begin{align*}
  V_D(\dev{\pi}{h}{a})-V_D(\pi)
  &=\sum_{s\in R_h(\pi)}\mu(s)\bigl(U_D(s,\dev{\pi}{h}{a})-U_D(s,\pi)\bigr)\\
  &=r_\pi(h)\bigl(Q_\pi(h,a)-Q_\pi(h,\pi(h))\bigr).
\end{align*}
If $r_\pi(h)>0$ and $\pi(h)$ is not causally optimal at $h$, the deviation raises the ex-ante value.
\end{proof}

Removing predictions about unreached information sets eliminates particular witnesses, such as
counterfactual mugging, and so changes the evidential part of the map in Theorem~\ref{thm:map16}.
It leaves the central Newcomb incompatibility intact.
